Table~\ref{tab:summary} summarises what worked and what did not; the subsections give the evidence. Unless stated otherwise, TVD values are in percentage points of total expenditure, CIs are 95\% cluster-bootstrap intervals over economies, and ``A16'' and ``A10'' denote the 16-function and ten-division COFOG panels (Section~\ref{sec:data}).

\begin{table}
\caption{Summary of the evidence: what worked and what did not (gains are TVD reductions in percentage points; Holm $p$ from Table~\ref{tab:holm}).}\label{tab:summary}
\centering\footnotesize
\renewcommand{\arraystretch}{1.12}
\begin{tabular}{@{}p{3.3cm}p{6.9cm}p{2.9cm}@{}}
\toprule
\textbf{Question} & \textbf{Evidence} & \textbf{Verdict}\\
\midrule
Learned correction beats last year (A16, nowcast)? & Gain \ResPairedGain{} (CI \ResPairedLo--\ResPairedHi) for learner seed 42; five-seed mean \AnnSeedMeanGain{}, CI excludes 0 in \AnnSeedNExcl{} of \AnnSeedN{} seeds & small, not robust\\
Same, wide panel A10? & stack \HistStackTen{} vs persistence \HistPersistTen{} (significantly worse) & no\\
Peers help without history, A16? & stack \FSStackSixteenGain{} (Holm $p=$ \FSStackSixteenHolm); 8-nearest \FSKnnSixteenGain{} (Holm $p=$ \FSKnnSixteenHolm) & weak\\
Peers help without history, A10? & stack \FSStackTenGain{} and 8-nearest \FSKnnTenGain{} (Holm $p=$ \FSStackTenHolm); XGBoost \FSXgbTenGain{} (Holm $p=$ \FSXgbTenHolm) & yes, modest ($\approx$8\% of a 15\,pp error)\\
Health and R\&D from indicators (WDI panel)? & $R^2$ 0.63 and 0.60 vs 0.31, 0.25 (peer mean); ridge 0.50, 0.53 & partly: mostly development level\\
Shock, institutional, WGI features; multi-year; econometric models? & none beats persistence or the baseline significantly & no\\
SHAP and LIME? & sign agreement \AgrSignMean, rank $\rho$ \AgrRhoMean; both recover 2 of 3 true drivers in a synthetic test & agree on direction, ambiguous for correlated features\\
Error by income group? & \ErrHigh{} (high) vs \ErrLow{} (low income, 3 economies) & worse for poorer economies\\
\bottomrule
\end{tabular}
\end{table}

\subsection{Forensic audit of the legacy benchmark}\label{sec:audit}
Before reporting new results we re-examine the legacy v1 benchmark, because its headline $R^2>0.99$ motivated this work. Three findings (Table~\ref{tab:audit}, Fig.~\ref{fig:audit}):
\begin{enumerate}
\item \textbf{One latent series, identical mix among eleven sectors.} Eleven of the fifteen sector columns are exact fixed multiples of one another (relative to agriculture: infrastructure $2.0\times$, social welfare $3.0\times$, transport $1.5\times$, energy $1.25\times$, housing $1.125\times$, rural development and industry $0.875\times$, water $0.75\times$, environment $0.625\times$, digital $0.5\times$; standard deviation of every ratio $\approx10^{-12}$). Among these eleven sectors all 775 country-years therefore share the same mix; only their total ($\approx16.7\%$ of GDP, CV 2.8\%, correlated $r=0.84$ with GDP growth) varies. There is no allocation pattern to learn in them; the other four sectors (defence, R\&D, education, health) do vary.
\item \textbf{GDP explains the $R^2$.} On these eleven sectors the one-parameter rule $k\times$GDP attains $R^2=\LegFormulaRConst$ versus $\LegXGBRConst$ for XGBoost on a random 80/20 split (USD scale), and a constant share of GDP gives $R^2=\LegConstRatioConst$ on the share scale. The mean over all fifteen sectors, $R^2=\LegXGBR$, is thus dominated by scale.
\item \textbf{Leakage and duplicated columns.} The legacy health-spending column is identical to the input feature of the same name (removing the feature changes $R^2$ from \LegHealthWith{} to \LegHealthWithout{} because other features still encode GDP), and \LegNIdent{} column pairs are identical, including all SDG ``budgets'' with their parent sector; the SDG-aligned model is therefore not an independent result.
\end{enumerate}
Only four legacy columns carry real information: education and R\&D equal the official WDI series (difference $<10^{-9}$ pp) and defence equals it to within 0.001\,pp (Pearson $r=0.998$), while healthcare differs from the official domestic government health series (median gap 2.4\,pp, $r=0.88$). The legacy \emph{features} match the World Bank values exactly for most indicators (inflation to within 0.4\% on 90\% of rows), which is why we keep and reuse them. With a temporal split on the legacy data XGBoost obtains $R^2=\LegTempXGB$ against \LegTempFormula{} for $k\times$GDP: again scale, not allocation. The legacy score therefore cannot support claims about learning budget-allocation patterns, and all conclusions below rest on official targets. The audit is a correction of our own earlier benchmark; we treat it as a methodological lesson on synthetic benchmarks, not as a main research result.

\begin{table*}
\caption{Forensic re-analysis of the legacy v1 benchmark (random 80/20 split as in the original notebooks). CV = coefficient of variation of sector/GDP across all 775 rows.}\label{tab:audit}
\centering\small
\renewcommand{\arraystretch}{1.1}
\resizebox{\linewidth}{!}{%
\begin{tabular}{rlcccccc}
\toprule
\# & Legacy sector & Mean \%GDP & CV (\%) & $R^2$ XGB (USD) & $R^2$ $k\!\times$GDP (USD) & $R^2$ XGB (share) & $R^2$ const. (share) \\
\midrule
1 & defence & 1.82 & 81.9 & 0.987 & 0.540 & 0.89 & $-$0.13 \\
2 & infrastructure & 2.47 & 2.8 & 0.999 & 1.000 & 0.95 & $-$0.00 \\
3 & research and development & 1.58 & 70.2 & 0.997 & 0.785 & 0.98 & $-$0.00 \\
4 & education & 5.43 & 79.4 & 0.978 & 0.823 & 0.76 & $-$0.03 \\
5 & healthcare & 7.31 & 44.3 & 0.987 & 0.709 & 0.98 & $-$0.00 \\
6 & agriculture & 1.24 & 2.8 & 0.999 & 1.000 & 0.91 & $-$0.00 \\
7 & rural development & 1.08 & 2.8 & 0.999 & 1.000 & 0.88 & $-$0.00 \\
8 & social welfare schemes & 3.71 & 2.8 & 0.999 & 1.000 & 0.96 & $-$0.00 \\
9 & energy & 1.54 & 2.8 & 0.999 & 1.000 & 0.93 & $-$0.00 \\
10 & transport & 1.85 & 2.8 & 0.999 & 1.000 & 0.93 & $-$0.00 \\
11 & housing and urban development & 1.39 & 2.8 & 0.999 & 1.000 & 0.92 & $-$0.00 \\
12 & water and sanitation & 0.93 & 2.8 & 0.999 & 1.000 & 0.85 & $-$0.00 \\
13 & environment and climate & 0.77 & 2.8 & 0.999 & 1.000 & 0.79 & $-$0.00 \\
14 & digital and IT & 0.62 & 2.8 & 0.999 & 1.000 & 0.70 & $-$0.00 \\
15 & industry and MSMEs & 1.08 & 2.8 & 0.999 & 1.000 & 0.88 & $-$0.00 \\
\bottomrule
\end{tabular}}
\end{table*}

\figAudit

\subsection{Annual update with history (Panel A16)}\label{sec:temporal}
Table~\ref{tab:A1} and Fig.~\ref{fig:main}a report the 2020--2024 test ($n=\ATest$ country-years; training: \ATrain{} country-years 2000--2019). \emph{Persistence is a strong baseline:} repeating last year's allocation requires moving only \ResTVDPersist\,pp of the budget (95\% CI \ResTVDPersistLo--\ResTVDPersistHi), compared with \ResTVDCountryMean\,pp for the country mean. This test window is hard: the same persistence rule misses by only 2.26\,pp in the 2017--19 validation window, so the COVID and energy-crisis years are about 60\% harder. Models that ignore history (features only) are clearly worse (6.3--9.2\,pp): a country's recent allocation carries far more information than its macro-indicators. The $\Delta$-learner (last year's shares plus a learned change) reduces TVD to \ResTVDHyb\,pp for the hybrid (CI \ResTVDHybLo--\ResTVDHybHi); for the pre-declared learner seed (42) the paired improvement over persistence is \ResPairedGain\,pp (95\% CI \ResPairedLo--\ResPairedHi, bootstrap over economies) and the hybrid is better on \ResRowsImproved\% of rows. This gain is \emph{small and not robust}: re-fitting the same hybrid with five learner seeds (42--46) gives gains of \AnnSeedMinGain--\AnnSeedMaxGain\,pp (mean \AnnSeedMeanGain\,pp) and a 95\% interval that excludes zero for only \AnnSeedNExcl{} of the \AnnSeedN{} seeds; XGBoost alone gains \AnnXgbSeedMeanGain\,pp on average and CatBoost provides most of the hybrid's gain. The individual learners (3.47--3.76\,pp) are statistically indistinguishable from the 50/50 hybrid, the log-ratio hybrid is indistinguishable from persistence (Section~\ref{sec:ext}), and this temporal comparison is not part of the Holm-corrected family. We therefore report the annual gain as a small positive point estimate that does not establish a reliable advantage over persistence. The learned change is informative about \emph{direction} for some functions (Fig.~\ref{fig:perfn}b: correlation 0.47 for education, 0.37 for general public services, 0.30 for health) but not for defence or R\&D (correlation $\approx0$).

\begin{table}
\caption{Temporal evaluation: train 2000--2019, test 2020--2024 ($n=178$ country-years, 36 economies); one-step-ahead with last year's allocation available. CIs: cluster bootstrap over economies. (MAE is omitted: with 17 components it equals TVD/8.5 exactly.)}\label{tab:A1}
\centering\small
\renewcommand{\arraystretch}{1.1}
\resizebox{\linewidth}{!}{%
\begin{tabular}{lccc}
\toprule
Method & TVD (pp) & 95\% CI & RMSE (pp) \\
\midrule
Persistence (last year) & 3.62 & [3.17, 4.08] & 0.923 \\
Country mean (train) & 6.96 & [6.16, 7.85] & 1.481 \\
Features-only: ridge regression & 9.22 & [8.37, 10.11] & 1.795 \\
Features-only: random forest & 6.45 & [5.72, 7.17] & 1.379 \\
Features-only: XGBoost & 6.64 & [5.94, 7.35] & 1.414 \\
Features-only: LightGBM & 6.34 & [5.68, 6.97] & 1.339 \\
Features-only: CatBoost & 6.44 & [5.78, 7.09] & 1.359 \\
Delta-learner: ridge regression & 3.76 & [3.31, 4.22] & 0.864 \\
Delta-learner: random forest & 3.47 & [3.04, 3.91] & 0.848 \\
Delta-learner: XGBoost & 3.59 & [3.15, 4.06] & 0.882 \\
Delta-learner: LightGBM & 3.63 & [3.17, 4.12] & 0.872 \\
Delta-learner: CatBoost & 3.50 & [3.08, 3.92] & 0.842 \\
Delta-learner hybrid (50\% XGBoost + 50\% CatBoost) & 3.48 & [3.06, 3.91] & 0.845 \\
\bottomrule
\end{tabular}}
\end{table}

\figMain

\textbf{COVID-19 and the energy crisis.} Fig.~\ref{fig:covid} compares predicted and observed one-year changes for 2020 and 2021. Models track the cross-country ranking of changes only weakly: for 2020 alone the correlation is \CovTwentyHealth\ for health and \CovTwentySocial\ for social protection, and for the pooled 2020--21 points in Fig.~\ref{fig:covid} it is \CovPooledHealth, \CovPooledSocial\ and \CovPooledEducation\ (health, social protection, education). They also miss the aggregate direction in 2020 (the mean health share rose by $+0.22$\,pp while the learned change was $-0.12$\,pp). Shock-year reallocations therefore remain the domain of human judgment.

\figPerfn
\figCovid

\subsection{Cold start: economies without history (Panel A16)}\label{sec:cold}\label{sec:selection}
In leave-country-out evaluation (Table~\ref{tab:A2}, Fig.~\ref{fig:main}b) the boosted hybrid reaches TVD \ResColdHyb\,pp (CI \ResColdHybLo--\ResColdHybHi) versus \ResColdPeer\,pp for the income-group peer mean and \ResColdGlobal\,pp for the global mean (two-seed means; the difference of the two means is $0.57$\,pp). The paired gain on fold-seed~0 only is \ResColdGain\,pp (95\% CI \ResColdGainLo\ to \ResColdGainHi), and the five-seed gain of single XGBoost is \FSXgbSixteenGain\,pp (CI \FSXgbSixteenLo\ to \FSXgbSixteenHi); neither is statistically significant, and ridge regression is worse than the peer mean. Composition across COFOG functions thus differs between economies for reasons the 25 macro-governance features capture only weakly (roughly one-tenth of the budget remains misallocated relative to the observed split). Note that the peer mean uses only four income classes and is the only simple baseline we ran; stronger simple baselines (region-by-income means, a GDP-per-capita regression) were not tested, and the ridge result suggests that linear baselines are also under-explored. Dropping any feature group raises error (Table~\ref{tab:ablation}; demographic/health indicators matter most) but all differences are within about two seed standard deviations. Richer feature specifications (region/income dummies, year-on-year changes) did \emph{not} help on a validation window and were not adopted (Table~\ref{tab:selection}).

\begin{table}
\caption{Cold-start evaluation: leave-country-out (6 folds, 2 fold seeds; TVD is the mean over both seeds), all 825 country-years, no history available.}\label{tab:A2}
\centering\small
\renewcommand{\arraystretch}{1.1}
\resizebox{\linewidth}{!}{%
\begin{tabular}{lccc}
\toprule
Method & TVD (pp) & 95\% CI & SD over fold seeds \\
\midrule
Hybrid (50\% XGBoost + 50\% CatBoost) & 10.77 & [9.60, 12.05] & 0.25 \\
random forest & 10.80 & [9.67, 12.07] & 0.12 \\
CatBoost & 10.83 & [9.65, 12.11] & 0.33 \\
XGBoost & 10.83 & [9.67, 12.12] & 0.16 \\
LightGBM & 10.97 & [9.79, 12.28] & 0.21 \\
Income-group peer mean & 11.34 & [10.33, 12.53] & 0.06 \\
Global mean & 11.35 & [10.35, 12.53] & 0.07 \\
ridge regression & 11.98 & [10.75, 13.26] & 0.18 \\
\bottomrule
\end{tabular}}
\end{table}


\subsection{Global evidence on four official functions (Panel B)}\label{sec:wdi}
The WDI panel (\NWdiCountries{} economies, 2000--2024) tests whether development indicators can predict the size of four official sector allocations in economies \emph{not} seen in training (Table~\ref{tab:B}, Fig.~\ref{fig:wdi}). In leave-country-out evaluation the hybrid explains $R^2=0.63$ of health spending (peer mean $0.31$) and $0.60$ of R\&D intensity (peer mean $0.25$); the gain is smaller for education ($0.30$ vs.\ $0.02$) and negligible for defence ($0.12$ vs.\ $-0.03$), whose level is driven by security context not captured by the features. A plain ridge regression already recovers most of this gain ($R^2$ 0.50 for health, 0.53 for R\&D), so the benefit of boosted trees over a linear model is modest ($+0.13$ and $+0.07$), and most of the advantage over the four-class peer mean reflects development level (GDP per capita). These targets are shares of GDP, whereas several features (central-government expense, tax revenue and debt, Table~\ref{tab:features}) describe the fiscal envelope from which the spending is financed; they are not transformations of the targets, but they are not independent of them either, and we did not re-run Panel B without them. Leave-\emph{region}-out is harder: health ($0.49$) and R\&D ($0.48$) remain predictable, education is weak ($0.16$), and defence and the shares-of-government-expenditure variants have \emph{negative} pooled $R^2$ under leave-region-out, so we caution against cross-region transfer for those functions. With history (temporal split), persistence is almost as good as the $\Delta$-hybrid ($R^2$ 0.93 vs.\ 0.94 for health, 0.86 vs.\ 0.87 for education, equal for defence and R\&D). Panel B is run with one fold seed, so no seed variability is reported.

\begin{table*}
\caption{WDI global panel, $R^2$. LCO = leave-country-out, LRO = leave-region-out: income-group peer mean vs.\ XGBoost+CatBoost hybrid. \emph{Temporal}: train $\le$2018, test 2019--2024, persistence vs.\ $\Delta$-hybrid with last-year value.}\label{tab:B}
\centering\small
\renewcommand{\arraystretch}{1.1}
\begin{adjustbox}{max width=\linewidth}
\begin{tabular}{lrcccccc}
\toprule
Official series & $n$ & LCO peer & LCO hybrid & LRO peer & LRO hybrid & Persist. & $\Delta$-hybrid \\
\midrule
Health (\%GDP) & 4547 & 0.31 & 0.63 & 0.24 & 0.49 & 0.93 & 0.94 \\
R\&D (\%GDP) & 2212 & 0.25 & 0.60 & 0.23 & 0.48 & 0.99 & 0.99 \\
Education (\%GDP) & 3365 & 0.02 & 0.30 & $-$0.10 & 0.16 & 0.86 & 0.87 \\
Defence (\%GDP) & 3814 & $-$0.02 & 0.12 & $-$0.25 & $-$0.09 & 0.82 & 0.82 \\
Health (\%gov. exp.) & 4547 & 0.28 & 0.43 & 0.25 & 0.14 & 0.89 & 0.90 \\
Education (\%gov. exp.) & 3304 & 0.06 & 0.13 & 0.02 & $-$0.13 & 0.78 & 0.78 \\
Defence (\%gov. exp.) & 3755 & 0.02 & 0.11 & $-$0.24 & $-$0.16 & 0.86 & 0.86 \\
\bottomrule
\end{tabular}\end{adjustbox}
\end{table*}

\figWdi

\subsection{Allocation engine (evaluated component)}\label{sec:engine}
The engine returns allocations that sum to 100\% to within $3\times10^{-14}$ and agree with an independent SLSQP solution to \NSlsqpDiff{} pp (maximum absolute difference, 100 rows). The width of the user-set bounds matters (Table~\ref{tab:bounds}, Fig.~\ref{fig:whatif}b): tight 1st--99th-percentile bounds clip 37\% of predictions and \emph{raise} error (TVD 3.65 vs.\ \EngTVDBefore\,pp without bounds) because \emph{observed} allocations frequently leave such narrow bands (43\% of test rows); bounds of $0.5\times$minimum to $1.5\times$maximum clip under 2\% of predictions with no accuracy loss. The projection is thus useful for enforcing the budget identity and for entering legal minima where they exist; it is not a mechanism to improve accuracy, and we do not count it as a methodological contribution.

\figWhatif

\subsection{Explainability}\label{sec:xai}
\paragraph{Global drivers.} Fig.~\ref{fig:shap} shows which features move each function's predicted share. Social protection is driven mainly by agriculture's value-added share (a proxy for development level) and the 65+ population share; education by population size, the 0--14 share and government expense; general public services by government expense, internet use and population; health by tax revenue, urbanisation, CO$_2$ per capita (a development proxy) and life expectancy; defence by trade openness. These are associations among correlated development indicators, not causal drivers, and the prominence of log population and government expense (the latter partly imputed, Table~\ref{tab:features}) shows that country size and fiscal-envelope variables enter the explanations.
\figShap
\paragraph{SHAP versus LIME (Table~\ref{tab:xai}, Fig.~\ref{fig:lime}).} Averaged over six functions, the rank correlation between $|\text{SHAP}|$ and $|\text{LIME}|$ is $\AgrRhoMean$, top-5 overlap (Jaccard) $\AgrJacMean$ and sign agreement on the top-5 features $\AgrSignMean$. The two methods therefore agree on \emph{direction} more than on \emph{ranking}; since a random sign has agreement 0.5 and most top features are monotone development proxies, part of the sign agreement is expected, and we did not compute a permuted-feature null. Individual cases can still disagree strongly: in the two case studies of Fig.~\ref{fig:lime} (Germany and Italy, health, 2020) the leading feature (tax revenue, respectively government expense, as \%GDP) has opposite signs under SHAP and LIME, because SHAP measures the contribution relative to the average prediction while LIME measures a local slope times the deviation from the mean, and attributions of correlated features can differ between the methods. A single explanation should therefore be read as a hypothesis. LIME local fits are modest ($R^2=\AgrLimeFit$ on average), and re-running LIME with another seed gives rank correlation $\AgrStabMean$ (0.75 for defence and general public services). In the deletion test, resetting the three most important SHAP features changes the prediction \AgrFaithMin--\AgrFaithMax$\times$ (mean \AgrFaithMean$\times$) as much as resetting three random features (six functions, 80 rows each). This shows that SHAP concentrates on features that move the model output; it is weak evidence of faithfulness because mean replacement creates off-manifold inputs and any importance ranking would be expected to pass, and we did not compare with permutation or gain importance. Neither method says that changing a feature would change a government's decision (Section~\ref{sec:disc}).
\begin{table*}
\caption{Quantitative evaluation of SHAP and LIME (80 random test rows per function). Faithfulness ratio = mean |prediction change| when the three most important SHAP features are reset to the training mean, divided by the same for three random features.}\label{tab:xai}
\centering\small
\renewcommand{\arraystretch}{1.1}
\begin{adjustbox}{max width=\linewidth}
\begin{tabular}{lcccccc}
\toprule
Function & Spearman $\rho$ & Top-5 Jaccard & Top-5 sign agr. & LIME local $R^2$ & LIME re-run $\rho$ & Faithfulness ratio \\
\midrule
Health & 0.58 & 0.41 & 0.82 & 0.47 & 0.92 & 3.1 \\
Social Protection & 0.64 & 0.52 & 0.90 & 0.59 & 0.92 & 3.4 \\
Defence & 0.52 & 0.43 & 0.84 & 0.56 & 0.75 & 3.4 \\
Education & 0.67 & 0.54 & 0.92 & 0.66 & 0.92 & 3.8 \\
Transport & 0.59 & 0.48 & 0.81 & 0.46 & 0.88 & 3.0 \\
General Public Services & 0.56 & 0.49 & 0.81 & 0.36 & 0.75 & 2.0 \\
\midrule
Mean & 0.59 & 0.48 & 0.85 & 0.52 & 0.86 & 3.1 \\
\bottomrule
\end{tabular}\end{adjustbox}
\end{table*}

\figLime
\paragraph{Synthetic ground truth.} In a synthetic test with known drivers (20 repetitions; one additive and one interaction term), TreeSHAP and LIME each recovered 2 of the 3 true drivers (precision@3 $=0.67$, identical in every repetition, i.e.\ a ceiling of two of three); when a near-copy of a driver was added, precision fell to 0.52 (TreeSHAP) and 0.43 (LIME) and attribution to the true driver dropped from 0.46 to 0.38 of the total. Correlated macro indicators therefore make single-feature attributions ambiguous, which is the setting of our real data.
\paragraph{What-if sensitivity.} Under a severe macro shock (GDP $-15\%$, inflation $+12$\,pp, debt $+35$\,pp; Fig.~\ref{fig:whatif}a) the cold-start model shifts shares by under 0.7\,pp for every function. General public services rises by 0.6\,pp (consistent with a larger debt-service burden) but social protection \emph{falls} by 0.7\,pp, contrary to the counter-cyclical expectation, because tree models cannot extrapolate beyond the support of the data and the shock moves several correlated features at once. Such scenarios are model sensitivities, not forecasts.

\subsection{Wider coverage (IMF data) and architecture v2.1}\label{sec:v21}
To test whether more and more diverse economies help, we added the IMF GFS expenditure-by-function dataset~\cite{imf_gfs} and built panel A10: the ten COFOG divisions for \NTenRows{} country-years of \NTenEcon{} economies, of which \NTenImf{} country-years (non-OECD economies such as India, Brazil, Indonesia) come from the IMF; on 723 overlapping country-years IMF and Eurostat agree to a median of at most 0.07\,pp (Pearson $r\ge0.984$ for all ten divisions). We also tested \emph{architecture v2.1}: history features, log-ratio targets and a non-negative least-squares \emph{stack} (Section~\ref{sec:method}).

\textbf{Annual update (Table~\ref{tab:Ehist}).} Neither richer history features nor log-ratio targets improve on the simple $\Delta$-learner. On panel A16 the stack (TVD \HistStackSixteen\,pp) is on par with the earlier hybrid and gains 0.11\,pp over persistence (paired, unrounded values; 95\% CI 0.03--0.20); on the wider panel A10 persistence is \HistPersistTen\,pp, the best learner \HistBestTen\,pp and the stack \HistStackTen\,pp: the stack is \emph{significantly worse} than persistence (paired $-0.20$\,pp, CI $-0.35$ to $-0.03$). For annual updates, persistence is the safer choice, and the simple $\Delta$-learner a marginal addition at best.

\textbf{Cold start (Tables~\ref{tab:Ecold} and~\ref{tab:holm}).} The stack has the lowest two-seed TVD on both panels. Averaged over five fold seeds, its paired gain over the income-group peer mean is \FSStackTenGain\,pp (95\% CI \FSStackTenLo--\FSStackTenHi, Holm-adjusted $p=$ \FSStackTenHolm) on A10 (TVD \FSTVDStackTen\ vs.\ \FSTVDPeerTen\,pp) and \FSStackSixteenGain\,pp (CI \FSStackSixteenLo--\FSStackSixteenHi, Holm $p=$ \FSStackSixteenHolm, not significant at 5\%) on A16 (\FSTVDStackSixteen\ vs.\ \FSTVDPeerSixteen\,pp). These are five-seed values and differ slightly from differences of the two-seed means in Table~\ref{tab:Ecold} (e.g.\ \ColdPeerTen{} $-$ \ColdStackTen{} $=1.24$ against 1.26). A single XGBoost gains \FSXgbTenGain\,pp (CI \FSXgbTenLo--\FSXgbTenHi) on A10 and \FSXgbSixteenGain\,pp on A16, neither surviving the Holm correction (the uncorrected A10 interval excludes zero), and the stack's advantage over XGBoost does not survive either. More, and more diverse, economies thus make peer learning useful on the wide panel, but the gain is about 1.3\,pp against an error of 15\,pp.

\begin{table}
\caption{Architecture v2.1, cold start (leave-country-out, 6 folds; TVD and CI are means over 2 fold seeds): peer mean, single learners (raw or log-ratio targets) and the nested-NNLS stack. Paired gains over the peer mean in the text are five-seed values (Table~\ref{tab:holm}), so they differ slightly from differences of the two-seed means shown here.}\label{tab:Ecold}
\centering\small
\renewcommand{\arraystretch}{1.1}
\begin{adjustbox}{max width=\linewidth}
\begin{tabular}{llcc}
\toprule
Panel & Method & TVD (pp) & 95\% CI \\
\midrule
A16 (36 economies) & Stack (nested NNLS) & 10.53 & [9.48, 11.67] \\
A16 (36 economies) & CatBoost (log-ratio targets) & 10.82 & [9.69, 12.04] \\
A16 (36 economies) & XGBoost (raw targets) & 10.83 & [9.67, 12.12] \\
A16 (36 economies) & XGBoost+CatBoost hybrid (log-ratio) & 10.86 & [9.72, 12.10] \\
A16 (36 economies) & XGBoost (log-ratio targets) & 11.07 & [9.91, 12.35] \\
A16 (36 economies) & Income-group peer mean & 11.34 & [10.33, 12.53] \\
A16 (36 economies) & Ridge & 11.98 & [10.75, 13.26] \\
\midrule
A10 (85 economies) & Stack (nested NNLS) & 13.77 & [12.30, 15.45] \\
A10 (85 economies) & XGBoost (raw targets) & 14.12 & [12.62, 15.78] \\
A10 (85 economies) & CatBoost (log-ratio targets) & 14.33 & [12.80, 16.13] \\
A10 (85 economies) & XGBoost+CatBoost hybrid (log-ratio) & 14.35 & [12.81, 16.13] \\
A10 (85 economies) & XGBoost (log-ratio targets) & 14.58 & [13.03, 16.34] \\
A10 (85 economies) & Ridge & 14.86 & [13.29, 16.70] \\
A10 (85 economies) & Income-group peer mean & 15.01 & [13.52, 16.80] \\
\bottomrule
\end{tabular}\end{adjustbox}
\end{table}

\begin{table}
\caption{Cold-start paired TVD gains averaged over five fold seeds, cluster bootstrap over economies (10\,000 draws), Holm correction over the family of eight comparisons. Positive = first method better.}\label{tab:holm}
\centering\small
\renewcommand{\arraystretch}{1.1}
\begin{adjustbox}{max width=\linewidth}
\begin{tabular}{llcccc}
\toprule
Panel & Comparison (cold start) & Gain (pp) & 95\% CI & Holm $p$ & Survives? \\
\midrule
A16 & Stack vs peer mean & 0.81 & [0.20, 1.34] & 0.060 & no \\
A16 & XGBoost vs peer mean & 0.47 & [-0.22, 1.13] & 0.171 & no \\
A16 & 8-nearest-economies vs peer mean & 0.66 & [0.19, 1.14] & 0.034 & yes \\
A16 & Stack vs XGBoost & 0.34 & [-0.03, 0.71] & 0.140 & no \\
A10 & Stack vs peer mean & 1.26 & [0.68, 1.84] & 0.0008 & yes \\
A10 & XGBoost vs peer mean & 0.89 & [0.19, 1.58] & 0.060 & no \\
A10 & 8-nearest-economies vs peer mean & 1.25 & [0.71, 1.78] & 0.0008 & yes \\
A10 & Stack vs XGBoost & 0.38 & [0.06, 0.69] & 0.060 & no \\
\bottomrule
\end{tabular}\end{adjustbox}
\end{table}


\textbf{Error by income group and applicability (Table~\ref{tab:income}).} Cold-start error rises steeply with income level: TVD \ErrHigh\,pp for high-income economies, 18.5 for upper-middle, \ErrLowerMid\ for lower-middle and \ErrLow\,pp for low-income economies, which have only nine and three training economies. An applicability score (Section~\ref{sec:method}) correlates with the error (Spearman $\rho=\NovRho$ over pooled rows), so the API reports it with every recommendation. This correlation is confounded by income group, which drives both novelty and error, and we did not compute a within-group partial correlation. Conclusions about low-income economies rest on three economies and are exploratory.

\begin{table}
\caption{Cold-start error (TVD, pp) by World Bank income group on the wide panel A10, with the mean applicability (novelty) score.}\label{tab:income}
\centering\small
\renewcommand{\arraystretch}{1.1}
\begin{adjustbox}{max width=\linewidth}
\begin{tabular}{lcccccc}
\toprule
Income group & Economies & Country-years & Novelty & Peer mean & Ridge & XGBoost \\
\midrule
High income & 50 & 1064 & 3.38 & 12.5 & 12.1 & 11.4 \\
Upper middle income & 23 & 371 & 4.05 & 18.3 & 19.4 & 18.5 \\
Lower middle income & 9 & 61 & 5.65 & 28.1 & 24.2 & 22.0 \\
Low income & 3 & 30 & 7.67 & 32.1 & 37.4 & 33.4 \\
\bottomrule
\end{tabular}\end{adjustbox}
\end{table}


\subsection{Total-budget module}\label{sec:total}
Shares become amounts only with a total. We therefore also model total general-government expenditure as a percentage of GDP (30 OECD/EU economies, 716 country-years with an observed total; the six OECD-only economies lack it) so that $\text{amount}_k=\text{share}_k\times\text{total}\times\text{GDP}$ (Table~\ref{tab:total}). In the temporal test a ridge $\Delta$-learner has a mean absolute error of \TotDeltaRidge\,pp of GDP against \TotPersist\,pp for persistence, because the 2020 shock moved total spending by several points of GDP and GDP growth is informative; for an economy without history the ridge model reaches \TotColdRidge\,pp against \TotColdMean\,pp for the group mean. Its feature set includes central-government expense as a percentage of GDP, which is a near-proxy for the target (correlation 0.75); in an internal audit check (not part of the released pipeline) dropping it raised the cold-start ridge error from 3.82 to 4.73\,pp, and dropping tax revenue, debt and expense together gave 5.97\,pp, worse than the group mean, so the module has no skill without those fiscal proxies. Because it is trained on OECD/EU economies only, its output for other economies is indicative; users may supply the total themselves.

\begin{table}
\caption{Total-budget module: general-government total expenditure as percent of GDP (30 OECD/EU economies, 716 country-years, Eurostat COFOG panel). MAE in percentage points of GDP. The income-group mean equals the global mean because all but one economy is high income.}\label{tab:total}
\centering\small
\renewcommand{\arraystretch}{1.1}
\resizebox{\linewidth}{!}{%
\begin{tabular}{llccc}
\toprule
Protocol & Method & $n$ & MAE (pp of GDP) & 95\% CI \\
\midrule
temporal 2020-24 & persistence & 150 & 2.87 & [2.57, 3.26] \\
temporal 2020-24 & country mean & 150 & 3.59 & [2.90, 4.49] \\
temporal 2020-24 & delta ridge & 150 & 1.83 & [1.62, 2.08] \\
temporal 2020-24 & delta xgb & 150 & 2.03 & [1.75, 2.38] \\
temporal 2020-24 & delta cat & 150 & 2.09 & [1.81, 2.43] \\
cold start leave-country-out & global mean & 716 & 5.63 & [4.47, 6.86] \\
cold start leave-country-out & income-group mean & 716 & 5.63 & [4.47, 6.86] \\
cold start leave-country-out & ridge & 716 & 3.82 & [3.13, 4.59] \\
cold start leave-country-out & xgb & 716 & 4.26 & [3.53, 5.02] \\
cold start leave-country-out & cat & 716 & 4.41 & [3.62, 5.24] \\
\bottomrule
\end{tabular}}
\end{table}


\subsection{Cold-start alternatives, intervals and fairness}\label{sec:ext}
\paragraph{Nearest economies.} The $k$-nearest-economies predictor ($k=8$, Section~\ref{sec:method}) is as accurate as the stack and XGBoost and is fully explainable: on A10 its five-seed TVD is 13.74\,pp against 14.11\,pp for XGBoost, and its gain over the peer mean is \FSKnnTenGain\,pp (CI \FSKnnTenLo--\FSKnnTenHi, Holm $p=$ \FSKnnTenHolm); on A16 the gain is \FSKnnSixteenGain\,pp (CI \FSKnnSixteenLo--\FSKnnSixteenHi, Holm $p=$ \FSKnnSixteenHolm). The value $k=8$ was fixed from a grid $\{3,5,8,12\}$ on out-of-fold data (where $k=12$ is marginally better) and the grid is not part of the Holm family, a mild forking-paths exposure. The API returns the five nearest economies with every recommendation.
\paragraph{Dirichlet regression and neural networks.} On the wide panel A10 the Dirichlet regression is as good as the stack (TVD 13.69\,pp against 13.77\,pp; 1.32\,pp better than the peer mean, CI 0.55--2.01) and the MLP reaches 13.85\,pp, whereas the transformer (14.14\,pp) matches XGBoost (14.08\,pp); on A16 none improves significantly on the peer mean or XGBoost. In the annual update no model beats persistence by a practically relevant margin (A10: best 4.58\,pp for a ridge $\Delta$-learner against 4.62\,pp; direct neural models without the persistence anchor are far worse, 6.55--6.91\,pp). A principled compositional model is therefore a strong, simple alternative to the stack for new economies, consistent with tree ensembles remaining strong on medium-sized tabular data~\cite{grinsztajn2022,shwartzziv2022}; our panels are smaller, so we read this only as a prior. (The stack, kNN and Dirichlet comparisons in this subsection and several below use fold-seed~0 or two-seed values, so the XGBoost cold-start TVD on A10 appears as 13.96 (fold-seed 0), 14.12 (two seeds) and 14.11\,pp (five seeds) in different tables; none changes a conclusion.)
\paragraph{Conformal intervals.} Split-conformal intervals calibrated on out-of-fold residuals of training economies reach near-nominal overall coverage (80\% $\to$ 80.1\%, 90\% $\to$ 90.3\%) but under-cover lower-income groups (low-income 47\%, lower-middle 61\% at the 80\% level), consistent with the lack of exchangeability discussed in Section~\ref{sec:method}; the website therefore warns users.
\paragraph{Hyper-parameter optimisation.} Optuna with nested grouped cross-validation on A10 gave no significant gain over the default settings (XGBoost 13.96 to 13.91\,pp, paired gain 0.05\,pp, CI $-0.19$ to 0.30; CatBoost 13.69 to 13.53\,pp, gain 0.17\,pp, CI $-0.09$ to 0.43; fold-seed 0), although the search was truncated by available compute; the defaults were kept.
\paragraph{Fairness by income group.} Cold-start error is largest for low-income economies (TVD \ErrLow\,pp against \ErrHigh\,pp for high-income). The max/min group ratio is 2.94 against 2.56 for the peer-mean baseline, so the learner does not close the gap for the lowest income group. The one detectable signed bias is an over-prediction of the social-protection share in low-income economies by 4.1\,pp. With three low-income economies, the bootstrap intervals for these quantities (23.7--40.9 for the error, 2.09--3.85 for the ratio, 2.5--7.0 for the bias) are not meaningful, and this audit is exploratory evidence of an equity concern rather than a formal fairness test.
\paragraph{Explaining the stacked blend.} Because the deployed recommendation is a blend, we also explained it directly with KernelSHAP (30 held-out country-years per division, panel A10) and compared it with TreeSHAP of the XGBoost member. The two explanations agree only moderately: the rank correlation of absolute attributions is 0.42 (health), 0.29 (education) and 0.37 (social protection), the top-5 overlap (Jaccard) is 0.32--0.39 and sign agreement on the top-5 features is 0.97, 0.71 and 0.77. KernelSHAP attributions on the blend change the blend 2.6, 1.3 and 2.8 times as much as random features when deleted. The explanation of the XGBoost member should therefore be read as indicative of, not identical to, the explanation of the blend.

\subsection{Negative results}\label{sec:neg}
\paragraph{Shock-aware features.} Growth, unemployment and debt changes, a recession dummy, a contemporaneous cross-country growth feature and a regime split did \emph{not} help: for 2020--21 on panel A16 the XGBoost shock-aware learner had TVD 4.16\,pp against 3.89\,pp for the baseline $\Delta$-learner (paired difference $-0.26$\,pp, 95\% CI $-0.41$ to $-0.11$) and 4.07\,pp for persistence; the CatBoost variant (4.01\,pp) was marginally below persistence but worse than its own baseline (3.73\,pp). In a global-financial-crisis hold-out (train up to 2007, test 2009--10) the XGBoost variant was likewise worse than persistence (4.07 vs.\ 3.43\,pp). Reallocation in crises remains a human judgment, in line with evidence on reallocation frictions~\cite{barnes2023}.
\paragraph{Multi-year projection.} Rolling origins (2012--19) and horizons of one to five years with recursive, frozen-macro and direct multi-horizon models: persistence was the best \emph{deployable} method at every horizon (A16: 2.92\,pp at $h=1$ and 5.63\,pp at $h=5$). A recursive projection with the \emph{true} future macro variables (an oracle) had TVD 2.80\,pp at $h=1$, a gain of 0.11\,pp over persistence that is not significant (CI $-0.13$ to 0.45), and was worse thereafter; frozen-macro projections were significantly worse. We do not offer multi-year projections.
\paragraph{Low-income coverage.} It cannot be repaired from the IMF source: its general-government COFOG file contains no values for 22 of 25 low-income and 35 of 47 lower-middle-income economies, and admitting central-government series with a relaxed consistency window adds only 13 country-years that are not comparable with general-government shares (mean TVD 7.5\,pp on the overlap). Inverse-frequency income weighting reduces low-income error only non-significantly (33.1 to 32.0\,pp).
\paragraph{Econometric panel baselines.} A random walk with country drift (TVD 3.73\,pp on A16, paired difference $-0.12$\,pp, CI $-0.15$ to $-0.09$), a pooled AR(1) on log-ratio shares with macro covariates (3.86\,pp) and the same with country fixed effects (3.88\,pp) are all significantly worse than repeating last year's allocation (3.62\,pp); on A10 the ordering is the same (persistence 4.62\,pp; 4.85--5.24\,pp for the econometric variants with penalty 1 and 4.63--5.0\,pp with penalty 10).
\paragraph{Robustness to fold seeds and multiple comparisons.} Table~\ref{tab:holm} gives five-seed gains with Holm correction over eight cold-start comparisons. Two claims survive at the 5\% level: the stack beats the income-group peer mean on A10 and the 8-nearest-economies predictor beats it on both panels. The stack on A16 ($p=0.060$), XGBoost over the peer mean ($p\ge0.06$) and the stack over XGBoost ($p\ge0.06$) do \emph{not} survive and should be read as suggestive. In the annual update, the log-ratio hybrid is statistically indistinguishable from persistence (raw $p=0.72$ on A16, $0.89$ on A10).
\paragraph{Institutional features.} Adding two keyless World Bank governance indicators (voice and accountability, political stability) and their two-year changes to the 25 features gave no significant gain: cold-start TVD 13.96 to 13.91\,pp (paired gain 0.04\,pp, CI $-0.12$ to 0.21) and, in the annual update, 4.74 to 4.71\,pp (gain 0.02\,pp, CI $-0.04$ to 0.08; 2020--21 only: 0.03\,pp). We then added three manually downloaded institutional sources: the IMF Fiscal Rules Dataset 1985--2024 (78\% of A10 rows covered), the Database of Political Institutions 2020 (election year and years to the next election; 77\%; ends in 2020) and the IMF MONA arrangements (programme active in the year; 100\%). Neither any single group nor all three together improved on the baseline. In cold start (three fold seeds) the paired gain (baseline minus extended TVD; negative means the extended model is \emph{worse}) was $-0.02$\,pp for fiscal rules (CI $-0.13$ to 0.09), $-0.02$\,pp for elections (CI $-0.09$ to 0.04), 0.00\,pp for IMF programmes (CI $-0.05$ to 0.06) and $-0.03$\,pp for all three (CI $-0.13$ to 0.08), i.e.\ TVD rose very slightly; in annual updates the gains were between $-0.016$ and $-0.006$\,pp with intervals that include zero (baseline 4.74\,pp). Institutional variables thus add no measurable predictive information beyond the economic, demographic and governance features already used, although coverage gaps (DPI ends in 2020) limit this conclusion.
